\clearpage % Rozdziały zaczynamy od nowej strony.
\section{Wymagania Projektowe}
\label{sec:wymagania}

Celem pracy jest opracowanie systemu analizy stylu jazdy kierowcy w oparciu o
telemetrię pozyskiwaną z magistrali CAN oraz model uczenia maszynowego, którego
wynik prezentowany jest kierowcy w czasie rzeczywistym. Skutkiem działania
systemu ma być poprawa ekonomiki jazdy oraz ograniczenie niekorzystnych nawyków
kierowcy, a w konsekwencji zmniejszenie zużycia paliwa i wolniejsze zużywanie
eksploatacyjnych podzespołów pojazdu. Realizacja tego założenia wymaga przejścia
przez kilka powiązanych ze sobą etapów.

Pierwszym krokiem w realizacji tego założenia jest analiza danych — z jednej strony
określenie, jakie parametry telemetryczne niosą informację o sposobie prowadzenia
pojazdu, takie jak prędkość pojazdu, prędkość obrotowa silnika oraz stopień
wciśnięcia pedału przyspieszenia i hamulca wraz z wyznaczanymi na ich podstawie
wielkościami pochodnymi, a z drugiej zbadanie, które z tych danych pojazd faktycznie
udostępnia na magistrali CAN oraz w jakiej postaci są one przesyłane. Na tym etapie
określa się również wymaganą częstotliwość oraz rozdzielczość akwizycji danych, tak
aby możliwe było wiarygodne odwzorowanie dynamiki jazdy.

Po ustaleniu, jakie dane mają być pozyskiwane i w jaki sposób, kolejnym krokiem jest
analiza wymagań stawianych platformie sprzętowej, pełniącej rolę jednostki
obliczeniowej oraz panelu kierowcy. Urządzenie musi dysponować interfejsem CAN
umożliwiającym odbiór danych z magistrali, wydajnością wystarczającą do lokalnego
wykonywania predykcji modelu uczenia maszynowego, wyświetlaczem, najlepiej
dotykowym, do prezentacji interfejsu kierowcy, slotem na kartę pamięci pozwalającym
zapisywać nagrania danych z magistrali, układem zegara czasu rzeczywistego oraz
łącznością bezprzewodową. Analiza dostępnych rozwiązań pozwala dobrać platformę
spełniającą wszystkie te wymagania w ramach jednego modułu.

Dysponując odpowiednią platformą, przechodzi się do zaprojektowania i
implementacji aplikacji odpowiedzialnej za gromadzenie danych z pojazdu. Obejmuje
to odbiór ramek z magistrali CAN w trybie wyłącznie nasłuchującym, który nie ingeruje
w komunikację pojazdu, ich parsowanie na podstawie identyfikatorów oraz walidację
polegającą na odrzucaniu wartości nieprawidłowych. Zarejestrowane dane są
znakowane czasem z wykorzystaniem układu zegara czasu rzeczywistego, którego
wskazania mogą być aktualizowane poprzez łączność bezprzewodową, a następnie
zapisywane na karcie pamięci w ustalonym formacie oraz przy stałej częstotliwości
próbkowania. Na tym etapie realizowane są również funkcje zapewniające
niezawodną pracę urządzenia w pojeździe, takie jak trwałe przechowywanie
parametrów konfiguracyjnych oraz mechanizmy oszczędzania energii, obejmujące
przechodzenie w stan uśpienia przy braku aktywności na magistrali i wybudzanie po
wykryciu ruchu.

Po zgromadzeniu danych kolejnym krokiem jest analiza modeli uczenia maszynowego
pod kątem ich przydatności do klasyfikacji stylu jazdy. Przeglądowi podlegają klasy
algorytmów nadające się do działania na urządzeniu wbudowanym, ze szczególnym
uwzględnieniem niskiego zużycia pamięci oraz krótkiego czasu predykcji. Na
podstawie tej analizy wybierane są modele kandydujące, które poddawane są
dalszym pracom.

Wybór modeli pozwala przejść do przygotowania danych uczących. Etap ten obejmuje
weryfikację zgromadzonych danych oraz usunięcie wartości błędnych, sprowadzenie
ich do stałej siatki próbkowania, a także odfiltrowanie fragmentów nieprzydatnych do
analizy, takich jak postój przy zgaszonym silniku czy jazda na biegu wstecznym. Na
podstawie przygotowanych danych wyznaczane są cechy opisujące charakter jazdy
w przesuwnych oknach czasowych, a poszczególnym oknom nadawane są etykiety
odpowiadające jeździe ekonomicznej, neutralnej oraz nieekonomicznej, zgodnie z
przyjętymi regułami opartymi na wartościach cech opisujących dynamikę jazdy.

Tak przygotowany zbiór służy do wytrenowania wybranych modeli, które są następnie
porównywane pod względem dokładności oraz stabilności działania, między innymi z
wykorzystaniem walidacji krzyżowej grupowanej po przejazdach. Wytrenowane
modele eksportowane są do postaci biblioteki w języku C, możliwej do wykorzystania
bezpośrednio w oprogramowaniu mikrokontrolera, z uwzględnieniem kwantyzacji
cech ograniczającej zużycie zasobów. Do konwersji modelu rozważane są dostępne
narzędzia, takie jak emlearn umożliwiający generowanie statycznego kodu w języku
C czy formaty pośrednie typu ONNX lub TensorFlow Lite.

Wyeksportowane modele osadzane są w oprogramowaniu panelu, gdzie predykcja
odbywa się lokalnie, na podstawie strumienia danych pochodzących z magistrali. Na
tym etapie przeprowadzana jest walidacja w warunkach rzeczywistych, podczas jazdy
pojazdem, pozwalająca ocenić, który z modeli radzi sobie lepiej, i wybrać rozwiązanie
docelowe. Na podstawie klasyfikacji kolejnych okien czasowych wyznaczany jest
bieżący wskaźnik oceny stylu jazdy, który jest podwyższany przy jeździe
ekonomicznej oraz obniżany w przypadku zachowań nieekonomicznych.

Następnym etapem jest zaprojektowanie i implementacja czytelnej oraz atrakcyjnej
wizualnie szaty graficznej interfejsu kierowcy. Wymaga to doboru biblioteki graficznej
umożliwiającej budowę interfejsu na wyświetlaczu wbudowanym oraz
zaprojektowania sposobu prezentacji wyników zgodnie z przyjętymi w celu pracy
zasadami nudgingu oraz gamifikacji, obejmującymi intuicyjną kolorystykę
sygnalizującą ocenę jazdy, bieżący wskaźnik punktowy oraz subtelne wskazówki
zamiast komunikatów o błędach. Istotne jest przy tym, aby interfejs nie rozpraszał
uwagi kierowcy i przyciągał ją jedynie wtedy, gdy jest to konieczne.

Ostatnim etapem prac są testy i walidacja całego rozwiązania, obejmujące weryfikację
poprawności akwizycji oraz parsowania danych z magistrali CAN, ocenę jakości
klasyfikacji modelu, a także testy działania systemu jako całości w warunkach
rzeczywistych. Testy obejmują zarówno weryfikację poszczególnych komponentów,
jak i pracę systemu jako całości, a ich wyniki podlegają dokładnej analizie i
dokumentowaniu, przy bieżącym usuwaniu wykrytych nieprawidłowości.

Najważniejszym etapem pracy będzie wykazanie, że połączenie telemetrii pojazdu z
metodami uczenia maszynowego oraz odpowiednio zaprojektowanym interfejsem
informacyjnym rzeczywiście pozwala skutecznie wpływać na styl jazdy kierowcy.
Wymaga to zdefiniowania mierzalnych kryteriów oceny, na podstawie których możliwe
będzie rozstrzygnięcie, w jakim stopniu przyjęte założenia znajdują potwierdzenie. Do
kryteriów tych należą skuteczność klasyfikacji stylu jazdy przez model, zdolność
interfejsu do wywoływania pożądanej zmiany zachowań kierowcy oraz przełożenie tej
zmiany na wymierne korzyści w zakresie ekonomii eksploatacji pojazdu, ograniczenia
emisji oraz bezpieczeństwa. Uzyskane wyniki podlegają analizie i interpretacji,
stanowiąc podstawę do oceny skuteczności opracowanego rozwiązania.

Na zakończenie projektu sporządzona zostanie dokumentacja techniczna opisująca
cały proces projektowania oraz implementacji rozwiązania. Dokumentacja ta powinna
obejmować opis architektury systemu i zastosowanych technologii, sposób
przygotowania danych oraz trenowania i wyboru modelu, schematy poszczególnych
komponentów, wyniki przeprowadzonych testów, a także wnioski oraz rekomendacje
dotyczące dalszego rozwoju systemu.